% !TEX root = ../hw01.tex
The identity $f(x)=x$ maps $(0,1)$ to itself, but it misses the
endpoints of $[0,1]$, as shown in Figure~\ref{fig:hw01:identity-map}.
We can make room for them by moving a countable sequence of points.

\begin{marginfigure}
  \centering
  % !TEX root = ../../problems/hw01.tex
\begin{tikzpicture}[x=2.9cm,y=2.9cm,>=Stealth,font=\footnotesize]
  \draw[->] (-0.12,0) -- (1.14,0) node[right] {$x$};
  \draw[->] (0,-0.12) -- (0,1.14) node[above] {$y$};
  \draw[thick] (0,0) -- (1,1);
  \draw[densely dotted,black!50] (1,0) -- (1,1) -- (0,1);
  \draw[fill=white] (0,0) circle[radius=1.8pt];
  \draw[fill=white] (1,1) circle[radius=1.8pt];
  \node[below left] at (0,0) {$0$};
  \node[below] at (1,0) {$1$};
  \node[left] at (0,1) {$1$};
  \node[above left] at (0.65,0.65) {$y=x$};
\end{tikzpicture}

  \caption{The identity on $(0,1)$ misses the target values $0$ and $1$.
    Both endpoints of the graph are excluded.}
  \label{fig:hw01:identity-map}
\end{marginfigure}

First assign
\[
  f\left(\frac12\right)=0,
  \qquad f\left(\frac13\right)=1.
\]
The values $1/2$ and $1/3$ must now be supplied by other inputs.
We use $1/4$ and $1/5$ for these values, and then continue shifting
the reciprocals to fill the resulting gaps. This gives
\[
  f\left(\frac14\right)=\frac12,\quad
  f\left(\frac15\right)=\frac13,\quad
  f\left(\frac16\right)=\frac14,\quad\ldots
\]
Figure~\ref{fig:hw01:reciprocal-shift} illustrates this shift.
We leave all other points fixed. Define a map $f$ from $(0,1)$
to $[0,1]$ by
\[
  f(x)=
  \begin{cases}
    0, & x=\frac12,\\[2pt]
    1, & x=\frac13,\\[2pt]
    \dfrac1{n-2}, & x=\frac1n,\quad n\in\N,\ n\ge4,\\[4pt]
    x, & \text{otherwise}.
  \end{cases}
\]

\begin{figure}[h]
  \centering
  % !TEX root = ../../problems/hw01.tex
\begin{tikzpicture}[x=0.92cm,y=1cm,>=Stealth,font=\small]
  \node[anchor=east] at (-0.35,1.2) {$x$};
  \node[anchor=east] at (-0.35,0) {$f(x)$};
  \foreach \pos/\source/\target in {
    0.2/{\frac12}/{0},
    1.55/{\frac13}/{1},
    2.9/{\frac14}/{\frac12},
    4.25/{\frac15}/{\frac13},
    5.6/{\frac16}/{\frac14},
    8.4/{\frac1n}/{\frac1{n-2}}} {
    \node at (\pos,1.2) {$\source$};
    \node at (\pos,0) {$\target$};
    \draw[->] (\pos,0.86) -- (\pos,0.35);
  }
  \node at (7,1.2) {$\cdots$};
  \node at (7,0) {$\cdots$};
  \node at (9.55,1.2) {$\cdots$};
  \node at (9.55,0) {$\cdots$};
\end{tikzpicture}

  \caption{The reciprocal sequence is shifted to accommodate both
    endpoints. Every point outside this sequence remains fixed.}
  \label{fig:hw01:reciprocal-shift}
\end{figure}

To check bijectivity, put
\[
  S=\left\{\frac1n\mid n\in\N,\ n\ge2\right\}.
\]
The cases in the definition are disjoint and cover $(0,1)$.
Every output lies in $[0,1]$, so the function is well-defined.

Among the inputs in $S$, only $1/2$ maps to $0$, and only $1/3$
maps to $1$. For each integer $m\ge2$, solving $1/(n-2)=1/m$
gives the unique input $1/(m+2)$. Thus $f$ maps $S$ bijectively
onto $S\cup\{0,1\}$.

On $(0,1)\setminus S$, the map is the identity, with image
$(0,1)\setminus S$. This image is disjoint from $S\cup\{0,1\}$,
so no fixed point can give a second preimage for a value already
supplied by the reciprocal sequence. Together the two images
cover $[0,1]$. Every target value therefore has exactly one
preimage, proving that $f$ is a bijection.
